\documentclass[12pt,a4paper,headings=small,version=first,dvips]{report}

\begin{document}

\title{NHLIB: a New Hazard LIBrary for OpenQuake}
\maketitle

\section*{Introduction}
Explain (IT) motivations for the creation of NHLIB:\\
- Improve PSHA calculation perfomances (reduce computation time, RAM usage)\\
- Introduce test coverage (required for any further software development)\\
- Facilitate integration with OpenQuake\\
Describe general structure of the document.\\

\section*{Source modeling}

Describe 4 source typologies and their usage (distributed seismicity, fault, subduction). Talk about rupture generators (ruptures are generated on the fly to avoid RAM consumption)

\subsection*{Area and Point source}
- explain area and point source parameters (figures included)\\
- show calculation examples (e.g.figure with rupture from area source with single nodal plane and hypo vs. area source with multiple(s) nodal planes and hypocenters)\\
- show code snippet (in the python shell)\\

\subsection*{Simple Fault source}
- explain simple fault source parameters (figures included)\\
- show calculation examples (e.g. figure with ruptures on a straight fault vs. non-straight fault, vertical vs. dipping...)\\
- show code snippet (in the python shell)\\

\subsection*{Complex Fault source}
- explain complex fault source parameters (figures included)\\
- show calculation examples (e.g. figure with ruptures on fault defined by single top and bottom edges, vs fault defined by multiple intermediate edges). Maybe take slab 1.0?\\
- show code snippet (in the python shell)\\

\section*{Ground Shaking Intensity modeling}
Describe GMPE and IPE (again code snipped to show how they can be initialized in a python shell)\\
Describe 'automatic testing', and table format\\
Describe vectorized calculations.\\

\section*{Hazard Curve calculator}
Describe Ned's equation, mention vectorized calculations.\\
Show example of hazard curve/map calculation with the different source typologies.

\section*{PEER Tests}
Report about implemented PEER tests. 

\section*{Performances}
Computation time for hazard Map calculation with single source.\\
RAM usage for area source (one with multiple nodal planes and finite ruptures)\\
Comparison with OpenSHA.

\section*{Conclusions}
Roadmap for future development.

\end{document}